\documentclass{article}
\usepackage{booktabs, graphicx, hyperref, fontspec}
\usepackage{sectsty}
\allsectionsfont{\sffamily}
\usepackage[margin=1in]{geometry}
\hypersetup{
  colorlinks = true,
  urlcolor = cyan,
 }
 \providecommand{\tightlist}{%
  \setlength{\itemsep}{0pt}\setlength{\parskip}{0pt}}
\newcommand*{\authorfont}{\fontfamily{phv}\selectfont}
\usepackage[]{Fira Sans}

\begin{document}

\sffamily

\centerline{\Huge Microeconomic Analysis}

\vspace{3 mm}

\centerline{\large Dr.~Ryan Safner}
\vspace{2 mm}
\centerline{\large \href{http://microF20.classes.ryansafner.com}{microF20.classes.ryansafner.com}}

\vspace{5 mm}

\begin{tabular}{@{}p{3.5in}p{3.5in}}           
\textbf{Course}: ECON 306 Fall 2020  & \textbf{Email:}  \href{mailto:safner@hood.edu}{\nolinkurl{safner@hood.edu}}\\
\textbf{Room}: ON ZOOM & \textbf{Office:}  Rosenstock 118\\
\textbf{Meets}: MW 2:00--3:15PM & \textbf{Hours:} TuTh 3:30--5:00PM ON ZOOM\\ 
\end{tabular}

\vspace{5 mm}

\hrule


\begin{quote}
``The curious task of economics is to demonstrate to men how little they
really know about what they imagine they can design.'' - F.A. Hayek,
1974 Economics Nobel Laureate
\end{quote}

\textbf{Microeconomic Analysis} is a collection of analytical tools for
modeling and understanding how decisions made by agents (individuals and
firms) in a context of market institutions lead to social coordination
and the creation of value. This toolkit is the analytical foundation for
exploring \emph{all} fields and applications of economics as a social
science, has many practical applications in business, finance, politics,
and insights into any realm of human endeavors. Economics, after all, is
the study of human action. At the intermediate level, these concepts are
made more rigorous through the use of mathematical models, but we will
first and foremost lay the theoretical foundations with economic
intuition. As such, this class assumes you have met the \textbf{required
prerequisite courses} by taking introductory economics\footnote{ECON 205
  \emph{and} ECON 206} and basic mathematics\footnote{MATH 112 \emph{or}
  ECMG 212}.

\hypertarget{hybrid-course-format}{%
\section*{Hybrid Course Format}\label{hybrid-course-format}}
\addcontentsline{toc}{section}{Hybrid Course Format}

This course is taught in a \textbf{hybrid} format, providing a mixture
of regular synchronous activity where we all can interact in real time,
with asynchronous material, which can be done remotely at your own pace.

\textbf{I will be holding all synchronous class sessions remotely (for
reasons I will make clear to you by the first day) on Zoom.} You can
attend these sessions on your computer or device from your dormitory or
home, and a classroom is available for you to use (socially-distanced,
and in masks), but I will not be in the classroom.

During the synchronous, scheduled times for the course (Monday/Wednesday
2:00 P.M.-3:15 P.M.), I will lecture on the material, hold in-class
discussions, and answer questions in real time \emph{on Zoom.}
Attendance to the live portion is strongly encouraged, but not required.

\textbf{Lecture slides, videos, and other synchronous materials will be
posted online by the end of the day in which the live session occurs.}

Assignments will always be submitted \emph{online} and due at regular
times (typically 11:59 PM Sunday) so that students unable to join in the
live sessions can complete them asynchronously.

Students are strongly encouraged to join the course
\href{https://hoodcollegeeconomics.slack.com}{Slack channel} to maintain
an active channel of communication, ask questions, and to build our
course community together. Official course-related announcements will
always come via Blackboard announcement and automatically sent to your
Hood email accounts.

\hypertarget{learning-in-a-time-of-coronavirus}{%
\subsection*{Learning in a Time of
Coronavirus}\label{learning-in-a-time-of-coronavirus}}
\addcontentsline{toc}{subsection}{Learning in a Time of Coronavirus}

Everything is awful right now. None of us signed up for this. None of us
are really okay, \textbf{we're all just pretending for everyone else.}

Many of you may be dealing with hardships at home and at work, and are
generally juggling many more problems than usual. Everyone's future
plans have been completely put on hold or cancelled to a large degree.
We all miss the sense of normalcy and human sense of community from
being isolated for so long.

For this unique semester, we are going to prioritize supporting each
other as human beings during this crazy era, and use simple, accessible
solutions that make sense for the most people, and above all, to be
flexible. I have designed the course to maintain some common structure
but be flexible to your varied needs. Please see the
\protect\hyperlink{policies-and-expectations}{policies and expectations
below}. I hope you use this course as an opportunity to escape the
boredom and insanity of social isolation, and to help keep interest in
understanding the world around us.

If you tell me you're having trouble, I will do whatever I can to help,
and not judge you or think less of you. I hope you will extend me the
same courtesy.

\hypertarget{course-objectives}{%
\section*{Course objectives}\label{course-objectives}}
\addcontentsline{toc}{section}{Course objectives}

This class will add rigor to your thinking about the world by empowering
you with the ability to model the choices and interactions of people and
organizations in their economic capacities -- as consumers, producers,
firms, voters, etc.

\textbf{By the end of this course,} you will:

\begin{enumerate}
\def\labelenumi{\arabic{enumi}.}
\tightlist
\item
  apply the models of microeconomics (constrained optimization and
  equilibrium) towards explaining real world behavior of individuals,
  firms, and governments
\item
  explore the effects of economic and political processes on market
  performance (competition, market prices, profits and losses, property
  rights, entrepreneurship, market power, market failures, public
  policy, government failures)
\item
  apply the economic way of thinking to real world issues in writing
\end{enumerate}

Given these objectives, this course fulfills all three of the
\href{https://www.hood.edu/academics/departments/george-b-delaplaine-jr-school-business/student-learning-outcomes}{learning
outcomes} for the Economics B.A. program in the George B. Delaplaine,
Jr.~School of Business:

\begin{itemize}
\tightlist
\item
  Use quantitative tools and techniques in the preparation,
  interpretation, analysis and presentation of data and information for
  problem solving and decision making {[}\ldots{]}
\item
  Apply economic reasoning and models to understand and analyze problems
  of public policy {[}\ldots{]}
\item
  Demonstrate effective oral and written communications skills for
  personal and professional success {[}\ldots{]}
\end{itemize}

\textbf{My standard disclaimer:} This class may challenge many of your
existing beliefs and conceptions about how the world works, and how it
should work. This is the most important and exciting part of a liberal
arts education. This does not mean that I want to make you to see
everything ``my way.'' In fact, if you come out of this class thinking
exactly like me, then I have probably failed you as a teacher. To the
best of my ability, I keep my opinions to myself unless relevant for
purposes of discussion, and I respect and invite you to reach your own
conclusions on all matters.

\textbf{Fair warning:} \textbf{Economics is hard.} \emph{This, in
particular, may be of the hardest courses that you will take, primarily
due to the mathematical content.} Using the economic way of thinking is
a skill, it is literally retraining your brain to interpret and analyze
the world in a novel way, and is not something that can be memorized. I
will do my best to make this class intuitive and helpful, if not
interesting. If at any point you find yourself struggling in this course
for any reason, please come see me. Do not suffer in silence! Coming to
see me for help does not diminish my view of you, in fact I will hold
you in \emph{higher} regard for understanding your own needs and taking
charge of your own learning. There are also a some fantastic resources
on campus, such as the
\href{http://www.hood.edu/campus-services/academic-services/index.html}{Center
for Academic Achievement and Retention (CAAR)} and the
\href{http://www.hood.edu/library/}{Beneficial-Hodson Library}.

See my
\href{http://microF20.classes.ryansafner.com/reference\#tips}{tips for
success in this course}.

\hypertarget{required-course-materials}{%
\section*{Required Course materials}\label{required-course-materials}}
\addcontentsline{toc}{section}{Required Course materials}

This course requires regular online internet access. If you know you
will be unable to access the internet regularly, please let me know and
we can make arrangements.

You can find all course materials at my \textbf{dedicated website} for
this course:
\href{https://microF20.classes.ryansafner.com}{microF20.classes.ryansafner.com}.
Links to the website are posted on our Blackboard course page. Please
familiarize yourself with the website, see that it contains this
\href{https://microF20.classes.ryansafner.com/syllabus/}{syllabus},
guides for your
\href{https://microF20.classes.ryansafner.com/reference/}{reference},
and our
\href{https://microF20.classes.ryansafner.com/schedule/}{schedule}. On
the schedule page, you can find each module with its own class page
(\textbf{start there!}) along with all related readings, lecture slides,
practice problems, and assignments.

My lecture slides will be shared with you, and serve as your primary
resource, but our main ``textbook'' below is \textbf{recommended} as the
next best resource and will be available from the campus bookstore. I
will discuss more about textbooks and materials in the first module.

\hypertarget{books}{%
\subsection*{Books}\label{books}}
\addcontentsline{toc}{subsection}{Books}

\begin{enumerate}
\def\labelenumi{\arabic{enumi}.}
\tightlist
\item
  Goolsbee, Austan, Steven Levitt, and Chad Syverson, 2012,
  \emph{Microeconomics} 2nd ed, USA: Worth Publishers, ISBN:
  978-1464187025
\end{enumerate}

You are welcome to purchase the book by other means (e.g.~Amazon,
half.com, etc). I have no financial stake in requiring you to purchase
this book. You are welcome to use previous version of the book, but
carefully verify the reading assignments, as the material may be
different across versions.

\hypertarget{assignments-and-grades}{%
\section*{Assignments and Grades}\label{assignments-and-grades}}
\addcontentsline{toc}{section}{Assignments and Grades}

Your final course grade is the weighted average of the following
assignments. You can find general descriptions for all the assignments
on the
\href{http://microF20.classes.ryansafner.com/assignments/}{assignments
page} and more specific information and examples on each assignment's
page on the
\href{http://microF20.classes.ryansafner.com/schedule/}{schedule page}.

\begin{center}

\begin{tabular}{lll}
\toprule
Frequency & Assignment & Weight\\
\midrule
1 & Opinion-Editorial & 20\%\\
7 & Homeworks & 20\% (using average HW grade)\\
3 & Exams & 20\% (each)\\
\bottomrule
\end{tabular}
\end{center}

Each assignment is graded on a 100 point scale. Letter-grade equivalents
are based on the following scale:

\begin{center}

\begin{tabular}{llll}
\toprule
Grade & Range & Grade1 & Range1\\
\midrule
A & 93–100\% & C & 73–76\%\\
A− & 90–92\% & C− & 70–72\%\\
B+ & 87–89\% & D+ & 67–69\%\\
B & 83–86\% & D & 63–66\%\\
B− & 80–82\% & D− & 60–62\%\\
\addlinespace
C+ & 77–79\% & F & < 60\%\\
\bottomrule
\end{tabular}
\end{center}

See also my
\href{https://ryansafner.shinyapps.io/306_grade_calculator/}{
\texttt{Grade\ Calculator}} app where you can calculate your overall
grade using existing assignment grades and forecast ``what if''
scenarios.

These grades are firm cutoffs, but I do of course round upwards
\((\geq\) 0.5) for final grades. A necessary reminder, as an academic, I
am not in the business of \emph{giving} out grades, I merely report the
grade that you \emph{earn}. I will not alter your grade unless you
provide a reasonable argument that I am in error (which does happen from
time to time).

\hypertarget{homeworks}{%
\subsection*{Homeworks}\label{homeworks}}
\addcontentsline{toc}{subsection}{Homeworks}

There will be roughly a homework assignment for every lesson. These
questions will help check your understanding and mastery of the
material, and will be a combination of quantitative, graph, and short
answer questions of examples. These will also be good indicators of what
to expect on the exams. You may collaborate with other students to work
on homeworks, but each person must turn in an individual assignment.
Homeworks are due one week from the class period where we finish a
lesson, and must be turned in at the beginning of class. Homeworks will
be graded by the T.A. as follows: 70\% of the grade for completion, and
30\% for one randomly selected question. This is to reward students for
putting in a full faith effort to try to reach an answer, even if not
every answer ends up being correct.

Homeworks are generally due by 11:59 PM Sunday by email. Please type
your answers to the following questions in a document and email a PDF
attachment to me.\footnote{In MS Word, or Pages, or most word processing
  software, File -\textgreater{} Save As -\textgreater{} PDF, or File
  -\textgreater{} Export -\textgreater{} PDF.} You may still handwrite
answers if you will be able to scan and email them, as a PDF, if they
are easily readable, but this is not preferred. It will be much easier
for me to grade, annotate, and return to you as a PDF, so please save me
some trouble. For the few questions that ask you to draw a graph, try to
do so on your computer (use MS Paint, the drawing tools in MS Word or MS
Powerpoint, plot points in MS Excel, drawing/notetaking apps, etc.), and
save it as an image to include on your homework document. Again, they
need not be perfect or to scale, just show that you understand the broad
idea. Being able to understand and sketch the graphs is still a very
important and useful skill! If all else fails, I will be lenient in
grading graph questions if you are unable to technologically include a
graph.

\hypertarget{exams}{%
\subsection*{Exams}\label{exams}}
\addcontentsline{toc}{subsection}{Exams}

There will be three exams (one at the end of each unit) including the
final exam. Each exam only covers the material in its associated unit
(i.e.~no exam is cumulative). These exams provide feedback both to you
and to me that ensures everyone is progressing on schedule and
apprehending the material. This is critical, as the rest of the course,
and indeed, any future economics course you take, will build off of this
foundation.

Exams will be released \textbf{on Blackboard} as a timed assignment. You
will have 2 hours once you open the exam on Blackboard. You will not
need the whole time, I have given you some extra to accomodate the
difficulties of taking an exam at home. Please pick a time to take it
where you know you will have 2 hours. You may close the exam page and
come back to it, but the timer will continue to run once the exam is
first opened.

\hypertarget{opinion-editorial}{%
\subsection*{Opinion-Editorial}\label{opinion-editorial}}
\addcontentsline{toc}{subsection}{Opinion-Editorial}

Economic fallacies have always been popular, and journalists,
politicians, and talking heads consistently engage in faulty economic
reasoning in print and on television. Your task, as a student of
economics, is to find some issue discussed in the past year, and write a
critique of media discussion of that issue. Alternatively, instead of
critiquing the reasoning or statements of others, you may write an
advocacy piece, where you propose some economic policy and argue to
persuade readers to endorse it. You will be graded both on the soundness
of your economic reasoning and the quality of your writing. I would be
happy to co-author an Op-Ed with anyone who writes an exemplary Op-Ed.
Students who successfully publish their work in a media outlet will earn
\emph{extra credit} on their Op-Ed grade. I will provide ample resources
and examples to help you accomplish this, you can find more on the
\href{}{assignment page}.

\hypertarget{no-extra-credit-is-available}{%
\subsection*{No extra credit is
available}\label{no-extra-credit-is-available}}
\addcontentsline{toc}{subsection}{No extra credit is available}

\hypertarget{policies-and-expectations}{%
\section*{Policies and Expectations}\label{policies-and-expectations}}
\addcontentsline{toc}{section}{Policies and Expectations}

This syllabus is a contract between you, the student, and me, your
instructor. It has been carefully and deliberately thought out\footnote{A
  syllabus can and will be used as a legal document for disputes tried
  at a court of law. Ask me how I know.}, and I will uphold my end of
the agreement and expect you to uphold yours.

In the language of game theory, this syllabus is my commitment device. I
am a very understanding person, and I know that exceptions to rules
often need to be made for students. However, to be \emph{fair} to
\emph{all} students the syllabus artificially constrains my ability to
make exceptions at a whim for anyone. This prevents clever students from
exploiting my congenial personality at everyone else's expense. Please
read and familiarize yourself with the course policies and expectations
of you. Chances are, if you have a question, it is answered herein.

\hypertarget{online-attendance-and-participation}{%
\subsection*{Online Attendance and
Participation}\label{online-attendance-and-participation}}
\addcontentsline{toc}{subsection}{Online Attendance and Participation}

This is a hybrid course with synchronous (live) and asynchronous (on
your own time) parts.

You are generally expected to join (online via Zoom) our
\textbf{synchronous} class sessions unless circumstances prevent you
from doing so. Day-to-day attendance is not graded per se, but I
strongly recommend you join in all live sessions in which you are able,
since we all can provide live feedback and I can answer questions and
address concerns as soon as they come up. You will also benefit from a
rigid schedule and shared community.

If you are unable to make a particular class, you generally do not need
to let me know. \textbf{The videos from all class sessions are posted on
Blackboard} so please review videos of classes you were unable to attend
live.

All assignmnents are able to be completed \textbf{asynchronously} during
the week, and are \textbf{generally due by 11:59PM Sunday each week} to
allow you flexibility in your hectic schedules.

\hypertarget{late-assignments}{%
\subsection*{Late Assignments}\label{late-assignments}}
\addcontentsline{toc}{subsection}{Late Assignments}

I will accept late assignments, but will subtract a specified amount of
points as a penalty. Even if it is the last week of the semester, I
encourage you to turn in late work: some points are better than no
points!

\textbf{Homeworks}: If you turn in a homework after it is due but before
it is graded or the answer key posted, I generally will not take off any
points. However, \textbf{if you turn in a homework \emph{after} the
answer key is posted, I will automatically deduct 20 points (so the
maximum grade you can earn on it is an 80).}

\textbf{Exams}: If you know that you will be unable to complete an
\emph{exam} as scheduled for a legitimate reason, please notify me at
least \emph{one week} in advance, and we will schedule a make-up exam
date. Failure to do so, including desperate attempts to make
arrangements only \emph{after} the exam will result in a grade of 0 and
little sympathy. I reserve the right to re-weight other assignments for
students who I believe are legitimately unable to complete a particular
assignment.

\textbf{Op-eds}: Starting at the deadline, I will take off 1 point for
every hour that your Op-ed is late.

\hypertarget{grading}{%
\subsection*{Grading}\label{grading}}
\addcontentsline{toc}{subsection}{Grading}

I will try my best to post grades on Blackboard's Grading Center and
return graded assignments to you within about one week of you turning
them in. There will be exceptions. Where applicable, I will post answer
keys once I know most homeworks are turned in (see Late Assignments
above for penalties). Blackboard's Grading Center is the place to look
for your most up-to-date grades. See also my
\href{https://ryansafner.shinyapps.io/306_grade_calculator/}{
\texttt{Grade\ Calculator}} app where you can calculate your overall
grade using existing assignment grades and forecast ``what if''
scenarios.

\hypertarget{communication-email-slack-and-virtual-office-hours}{%
\subsubsection*{Communication: Email, Slack, and Virtual Office
Hours}\label{communication-email-slack-and-virtual-office-hours}}
\addcontentsline{toc}{subsubsection}{Communication: Email, Slack, and
Virtual Office Hours}

Students must regularly monitor their \textbf{Hood email accounts} to
receive important college information, including messages related to
this class. Email through the Blackboard system is my main method of
communicating announcements and deadlines regarding your assignments.
\textbf{Please do not reply to any automated Blackboard emails - I may
not recieve it!}. My Hood email (\texttt{safner@hood.edu}) is the best
means of contacting me. I will do my best to respond within 24 hours. If
I do not reply within 48 hours, do not take it personally, and
\emph{feel free to send a follow up email} in the very likely event that
I genuinely did not see your original message.

Our \href{https://hoodcollegeeconomics.slack.com}{slack channel} is
available to all students and faculty in Economics and Business. I have
invited all of my classes and advisees. It will not be extended to
non-Business/Economics students or faculty. All users must use their
\textbf{hood emails} and \textbf{true first and last names}. Each course
has its own channel, exclusive for verified students in the course, and
myself, by my invite only. As a third party platform, you agree to its
Terms of Service. I have created this space as a way to stay connected,
to help one another, and to foster community. Behaviors such as posting
inappropriate content, harassing others, or engaging in academic
dishonesty, to be determined solely at my discretion, will result in one
warning, the content will be deleted, and subsequent behavior will
result in a ban.

I will host general \textbf{``office hours''} on Zoom. You can join in
with video, audio, and/or chat, whichever you feel comfortable with. Of
course, if you are not available during those times, we can schedule our
own time if you prefer this method over email or Slack. If you want to
go over material from class, please have \emph{specific} questions you
want help with. I am not in the business of giving private lectures
(particularly if you missed class without a valid excuse).

Watch the excellent and accurate video
\href{https://vimeo.com/270014784}{explaining office hours} (on website
syllabus page).

\hypertarget{netiquette}{%
\subsection*{Netiquette}\label{netiquette}}
\addcontentsline{toc}{subsection}{Netiquette}

When using Zoom and Slack, please follow appropriate internet etiquette
(``Netiquette''). Written communications, like blog posts or use of the
Zoom chat, lacks important nonverbal cues (such as body language, tone
of voice, sarcasm, etc).

Above all else, please respect one another and think/reread carefully
about how others may see your post before you submit a comment. You are
expected to disagree and have different opinions, this is inherently
valuable in a discussion. Please be civil and constructive in responding
to others' comments: writing \emph{``have you considered `X'?''} is a
lot more helpful to all involved than just writing \emph{``well you're
just wrong.''}

Posting content that is wilfully incindiary, illegal, or that
constitutes academic dishonesty (such as plagarism) will automatically
earn a grade of 0 and may be elevated to other authorities on campus.

When using the chat function on Zoom or public Slack channels, please
treat it as official course communications, even though I may not be
grading it. It may be a quick and informal tool - don't feel you need to
worry about spelling or perfect grammar - but please try to avoid
\emph{too} informal ``text-speak'' (i.e.~say ``That's good for you''
instead of ``thas good 4 u'').

\hypertarget{privacy}{%
\subsection*{Privacy}\label{privacy}}
\addcontentsline{toc}{subsection}{Privacy}

\href{https://www.execvision.io/blog/maryland-call-recording-laws/}{Maryland
law}
\href{https://law.justia.com/codes/maryland/2005/gcj/10-402.html}{requires}
all parties consent for a conversation or meeting to be recorded. If you
join in, and certainly if you participate, \textbf{you are consenting to
be recorded.} However, as described below, videos are \emph{not
accessible} beyond our class.

Live lectures are recorded on Zoom and posted to Blackboard via Panopto,
a secure course management system for video. Among other nice features
(such as multiple video screens, close captioning, and time-stamped
search functions!), Panopto is authenticated via your Blackboard
credentials, ensuring that \emph{our course videos are not accessible to
the open internet.}

For the privacy of your peers, and to foster an environment of trust and
academic freedom to explore ideas, \textbf{do not record our course
lectures or discussions.} You are already getting my official copies.

The
\href{https://www2.ed.gov/policy/gen/guid/fpco/ferpa/index.html}{Family
Educational Rights and Privacy Act} prevents me from disclosing or
discussing any student information, including grades and records about
student performance. If the student is at least 18 years of age,
\emph{parents (or spouses) do not have a right to obtain this
information}, except with consent by the student.

Many of you may be tuning in remotely, living with parents, and may have
occasional interruptions due to sharing a space. This is normal and
fine, but know that I will protect your privacy and not discuss your
performance when parents (or anyone other than you, for that metter) are
present, without your explicit consent.

\hypertarget{enrollment}{%
\subsection*{Enrollment}\label{enrollment}}
\addcontentsline{toc}{subsection}{Enrollment}

Students are responsible for verifying their enrollment in this class.
The last day to add or drop this class with no penalty is
\textbf{Thursday, August 27}. Be aware of
\href{https://www.hood.edu/offices-services/registrars-office/academic-calendar}{important
dates}.

\hypertarget{honor-code}{%
\subsection*{Honor Code}\label{honor-code}}
\addcontentsline{toc}{subsection}{Honor Code}

Hood College has an Academic Honor Code which requires all members of
this community to maintain the highest standards of academic honesty and
integrity. Cheating, plagiarism, lying, and stealing are all prohibited.
All violations of the Honor Code are taken seriously, will be reported
to appropriate authority, and may result in severe penalties, including
expulsion from the college. See
\href{http://hood.smartcatalogiq.com/en/2016-2017/Catalog/The-Spirit-of-Hood/The-Academic-Honor-Code-and-Code-of-Conduct}{here}
for more detailed information.

\hypertarget{van-halen-and-mms}{%
\subsection*{Van Halen and M\&Ms}\label{van-halen-and-mms}}
\addcontentsline{toc}{subsection}{Van Halen and M\&Ms}

When you have completed reading the syllabus, email me a picture of the
band Van Halen and a picture of a bowl of M\&Ms.~If you do this
\emph{before} the date of the first exam, you will get bonus points on
the exam.\footnote{If 75-100\% of the class does this, you each get 2
  points. If 50-75\% of the class does this, you each get 4 points. If
  25-50\% of the class does this, you each get 6 points. If 0-25\% of
  the class does this, you each get 8 points.} Yes, this is real.

\hypertarget{accessibility-equity-and-accommodations}{%
\subsection*{Accessibility, Equity, and
Accommodations}\label{accessibility-equity-and-accommodations}}
\addcontentsline{toc}{subsection}{Accessibility, Equity, and
Accommodations}

College courses can, and should, be challenging and bring you out of
your comfort zone in a safe and equitable environment. If, however, you
feel at any point in the semester that certain assignments or aspects of
the course will be disproportionately uncomfortable or burdensome for
you due to any factor beyond your control, please come see me or email
me. I am a very understanding person and am happy to work out a solution
together. I reserve the right to modify and reweight assignments at my
sole discretion for students that I belive would legitimately be at a
disadvantage, through no fault of their own, to complete them as
described.

If you are unable to afford required textbooks or other resources for
any reason, come see me and we can find a solution that works for you.

This course is intended to be accessible for all students, including
those with mental, physical, or cognitive disabilities, illness,
injuries, impairments, or any other condition that tends to negatively
affect one's equal access to education. If at any point in the term, you
find yourself not able to fully access the space, content, and
experience of this course, you are welcome to contact me to discuss your
specific needs. I also encourage you to contact the
\href{https://www.hood.edu/academics/josephine-steiner-center-academic-achievement-retention/accessibility-services}{Office
of Accessibility Services} (301-696-3421). If you have a diagnosis or
history of accommodations in high school or previous postsecondary
institutions, Accessibility Services can help you document your needs
and create an accommodation plan. By making a plan through Accessibility
Services, you can ensure appropriate accommodations without disclosing
your condition or diagnosis to course instructors.

\hypertarget{acbsp-information}{%
\subsection*{ACBSP Information}\label{acbsp-information}}
\addcontentsline{toc}{subsection}{ACBSP Information}

\begin{figure}[h!]
\centering
\includegraphics[]{../images/acbsp.png}
\end{figure}

Hood College is an accredited member of the Accreditation Council for
Business Schools and Programs (ACBSP), an organization devoted to
enhancing business education. In receiving and maintaining this
accreditation, the faculty has made a commitment to the continuous
improvement, innovation, and scholarship of the Department of Economics
and Management. For you, this means that your educational experience
undergoes ongoing validation to ensure it meets the most rigorous
international standards of business education. Only a select group of
institutions have received this status and it is an attribute of Hood
for which you should take great pride.

Pragmatically, our accreditation means that we will engage in the
ongoing use of measures, both quantitative and qualitative in nature, to
assess the performance of our students and the program. We ask that you
take very seriously the surveys and other measurement devices we will
use -- your best work and honest response will help us best assess and
improve our program.

\hypertarget{tentative-schedule}{%
\section*{Tentative Schedule}\label{tentative-schedule}}
\addcontentsline{toc}{section}{Tentative Schedule}

Below is a rough sketch of the weekly schedule we will aim to follow
this semester. Each module should take approximately one class meeting.

\textbf{You can find a full schedule} with much more details, including
the readings, appendices, and other further resources for each class
meeting on the
\href{http://microF20.classes.ryansafner.com/schedule/}{course website's
schedule page}.

\begin{center}
\small

\begin{tabular}{llll}
\toprule
Week & Topics & Readings & Assignments\\
\midrule
8/16-8/22 & Tools of Microeconomics; Scarcity, Choice and Cost & None & Survey\\
8/23-8/29 & Budget Constraint; Preferences & Goolsbee Ch. 4 & HW 1\\
8/30-9/5 & Solving the Consumer's Problem; Building a Demand Function & Goolsbee Ch. 5 & HW 2\\
9/6-9/12 & Income \& Substitution Effects; Price Elasticity & Goolsbee Ch. 5 & HW 3, Exam 1\\
9/13-9/19 & Theory of the Firm; Short Run \& Long Run & Goolsbee Ch. 6 & None\\
\addlinespace
9/20-9/26 & Cost Minimization; Costs of Production & Goolsbee Ch. 7 & HW 4\\
9/27-10/3 & Short Run Profit Maximization; Long Run Industry Equilibrium & Goolsbee Ch. 8 & HW 5, Exam 2\\
10/4-10/10 & Supply and Demand Model; Market Competition \& Surpluses & Goolsbee Ch. 2 & None\\
10/11-10/17 & Social Functions of Market Prices; Monopoly & Goolsbee Ch. 9 & HW 6\\
10/18-10/24 & Consequences \& Sources of Market Power; Pricing with Market Power & Goolsbee Ch. 9 & HW 7\\
\addlinespace
10/25-10/31 & Monopolistic Competition; Oligopoly & Goolsbee Ch. 10 & None\\
11/1-11/7 & Game Theory & Goolsbee Ch. 11 & None\\
11/8-11/14 & Factor Markets; Market Failures & Goolsbee Ch. 13 & None\\
11/15-11/21 & Market Failures & Goolsbee Ch. 14 & None\\
11/22-11/28 & Review & None & Final Exam\\
\bottomrule
\end{tabular}
\end{center}

\end{document}